\begin{slikaDesno}{fig/R.pdf}
\PID 
Приликом мерења непознате отпорности $R$, измерени су напон на њему 
$V = (1,03 \pm 0,01)\unit{V}$, и струја $I = (5,7 \pm 0,2)\unit{mA}$,
са одговарајућим апсолутним грешкама, 
према референтним смеровима са слике. Мерења напона и струје 
могу се сматрати некорелисанима.
\begin{enumerate}[label=(\alph*)]
\item Проценити вредност мерења 
$R$, одговарајућу апсолутну грешку $\Delta R$, и 
релативну грешку $\updelta R$ тог мерења. 
\item Поновљеним мерењем, са истом мерном поставком, добијена је средња 
вредност напона $\overline V = 1,00\unit{V}$, и средња вредност мерења 
струје $\overline I = 5,6\unit{mA}$. Обављено је $n = 100$ мерења. Одредити нове вредности 
$R$, $\Delta R$, и $\updelta R$. Сматрати да су $u_V = \Delta V$, $u_I = \Delta I$, и 
$u_R = \Delta R$.
\end{enumerate}
\end{slikaDesno}

\REZULTAT

(а) Измерена отпорност је $R = (180 \pm 8)\unit{\Omega}$, а релативна грешка је $\updelta R = 4,5\%$. 

(б) Измерена отпорност је $R = (178,6 \pm 0,8)\unit{\Omega}$.